\section{The step-(ii) density law and the off-grid frontier}
\label{sec:frontier}

This section reports the two newest components of the program: an
unconditional local bound on the sieve masses together with the exact
decay law it implies, and the current state of the off-grid horizon.

\subsection{An unconditional local bound}

Fix an aligned class with parameters $(\varepsilon,f,q_1,N)$ and write
$F(k) = P(\varepsilon f (q_1 + kN))$ for its degree-$8$ class
polynomial, and $u(k) = 2\alpha\varepsilon f(q_1+kN)$ for the unit whose
Legendre symbol decides bad signature. For an odd prime $p$ let $m_p$
denote the Haar mass in $\Zp$ of the set of $k$ with $v_p(F(k))$ odd and
$\leg{u(k)}{p} = -1$: the local mass of the emergent obstruction.

\begin{theorem}[Local sieve bound; PROVED]
\label{thm:localbound}
For every aligned class of the grid,
\[
   m_3 = m_5 = m_7 = 0,
   \qquad\text{and}\qquad
   m_p \;\le\; \frac{8}{p+1} \;<\; 1
   \quad\text{for all } p \ge 11 ,
\]
the bound holding for singular roots as well, counted with
multiplicity.
% src: data/l17_stepii.jsonl (meta.labels.theorem_replays = PROVED)
\end{theorem}

The vanishing at $p \in \{3,5,7\}$ is exactly the frozen-alignment
condition of the class certificate: those places are controlled, and
their symbols are pinned to $+1$ by construction. For $p \ge 11$ the
bound is the root count of a degree-$8$ polynomial: at most eight
residues can be roots, each simple root contributes Haar mass
$1/(p+1)$ for odd valuation, and a singular root contributes at most
its multiplicity times that mass, which is again bounded by the total
degree. The bound was verified independently on every stored local
pair: $102{,}439$ pairs with $p \le 10^4$ across the $293$ classes,
zero violations.
% src: data/l17_stepii.jsonl (verification.positive_local_pairs_checked_p_le_10000 = 102439, bound_violations = 0)

Theorem~\ref{thm:localbound} settles a natural hope in the negative,
and this is the important structural point.

\begin{corollary}[No uniform positive mass; PROVED for the current classes]
\label{cor:nouniform}
For each of the $293$ classes the splitting condition governing the
bad-signature event has Chebotarev exponent $c = 1/2$: the class
polynomial is irreducible with a replayed degree-$8$ certificate, and a
replayed witness shows the bad unit is a nonsquare at a simple bad
root. Consequently the truncated clean product obeys
\[
   \prod_{p \le X} (1 - m_p) \;\asymp\; (\log X)^{-1/2},
\]
so no cutoff-uniform positive lower bound on the clean mass exists:
the small-layer clean probability decays, slowly, to zero.
% src: data/l17_stepii.jsonl (verdict; serialized_witnesses_replayed = 293)
\end{corollary}

The measured exponent agrees: the root-harmonic estimate gives
$c \approx 0.4778$, the product estimate $c \approx 0.4801$, and the
same-member sieve between $10^4$ and $10^5$ gives $c \approx 0.4590$
(EVIDENCE).
% src: data/l17_stepii.jsonl (evidence.root_harmonic_c_mean, product_c_mean, same_member_sieve_c_10000_100000)

The consequence for hypothesis $H$ is a sharpening, not a defeat. The
number of admissible prime members up to height $X$ in a class grows
like $X/(\log X)^{2}$ by the Dirichlet--Bateman--Horn heuristic
\cite{batemanhorn}, while
each carries clean probability $\asymp (\log X)^{-1/2}$; the expected
count of emergent-free members therefore still diverges, at rate
$X/(\log X)^{5/2}$. What Corollary~\ref{cor:nouniform} rules out is the
clean shortcut --- a uniform positive constant --- and what it leaves is
a divergence statement of exactly the shape the measured counting law
already exhibits (first emergent-free member at $k=0$ in $207$ of $293$
classes).

\subsection{The off-grid horizon}

The grid is $w$ prime, $w<100$.  Off-grid closures now exist for
\[
 w\in\{101,103,107,109,113,127,131,137,139,149,151,157,163,167,173,179\},
 \qquad u=-1,
\]
together with $[131,[5,1]]$, each certified by an aligned class and an exact
ladder verdict \emph{zero}.  The routes are the L11 family, the canonical
fixed-factor escape, and the non-canonical $a=7$, $\tau=0$ family.  The
last searched open cell, $[179,[-1,1]]$, is now closed, so every off-grid
cell tried to date has a replayed member.  This is a finite-horizon
statement only.
% src: data/l17_horizon*.jsonl; data/l18_horizon_sweep_a.jsonl
% src: data/l18_horizon_sweep_b.jsonl; data/l18_route131.jsonl
% src: data/l19_cell179.jsonl (closure and summary rows)

Two negative or delimiting findings carry structural content.

\begin{theorem}[The canonical $5$-wall; PROVED]
\label{thm:fivewall}
On the canonical escape route ($z=-w$, $a=1$, $\tau=3/5$, $f=w$), write
$e = v_5(D)$. Then $v_5(x_0) = -3-2e$ is odd and $v_5(d_0)=0$, so
\[
   \mathrm{Hilb}_5 \;=\; -\leg{w}{5}\leg{q_1}{5},
   \qquad
   \mathrm{Hilb}_{q_1} \;=\; \leg{5}{q_1} ,
\]
and both equal $+1$ for some admissible $q_1$ if and only if
$\leg{w}{5} = -1$. Since the L11 route requires a root of
$1+4r^2 \equiv 0 \pmod w$, i.e.\ $w \equiv 1 \pmod 4$, the canonical
protocol is obstructed exactly for
\[
   w \equiv 3 \pmod 4 \ \text{ and } \ \leg{w}{5} = +1
   \qquad\Longleftrightarrow\qquad
   w \equiv 11, 19 \pmod{20}.
\]
% src: data/l18_fivewall.jsonl (proved identity + 53-prime audit, 0 mismatches)
\end{theorem}

The criterion was audited on all $53$ primes $w \in [101,400]$ with both
character branches of $q_1$ covered: $26$ cells have the L11 route, $13$
the canonical escape route, $14$ are obstructed, and there were zero
identity or classification mismatches.
% src: data/l18_fivewall.jsonl (classification_counts, identity_mismatches = 0)

The wall is a statement about the canonical parameter, not about the cells.
In fact it is always breakable on the class side.

\begin{theorem}[Every canonical $5$-wall is breakable; PROVED]
\label{prop:break}
Let $w\equiv11,19\pmod{20}$ be prime and set $z=-w$, $f=w$.  There is an
odd integer $a$ such that the $\tau=0$ fixed-factor family has an aligned
class for infinitely many positive primes $q_1$.  This theorem concerns
class existence; it does not assert an emergent-free member.  If $a$ is
fixed to $7$, the class-side break exists exactly when
$\leg{w}{197}=-1$.
\end{theorem}
% src: data/l19_tauzero.jsonl (criterion and always-breakable theorem rows)

\begin{proof}
Because $w\equiv3\pmod4$, the character sum
$\sum_{r\bmod w}\leg{1+4r^2}{w}=-1$ has no zero terms.  Exactly
$(w+1)/2$ residues therefore give character $-1$; choose one for $a\bmod
w$.  By CRT also impose
$a\equiv0\pmod{\operatorname{rad}(w^3+2)}$ and choose an odd lift.  For
$A=1+4a^2$, $Z=-w^3$, and $D=1-Z-a^2Z^2$, the identity
\[
 4D\equiv(Z-2)^2\pmod p\qquad(p\mid A)
\]
then gives $\gcd(A,D)=1$ while preserving $\leg{A}{w}=-1$.

At every prime occurring to odd order in $A$, prescribe
$\leg{-2\varepsilon wq_1}{p}=+1$; impose the exact $2$-adic class
$\varepsilon wq_1\equiv a\pmod4$; and make
$-2A\varepsilon wq_1$ a square at each of $3,5,7$ not dividing $A$.
These independent local conditions define a reduced CRT class, so
Dirichlet gives infinitely many primes $q_1$ in it.  After finitely many
exclusions and a real-sign threshold, every frozen, moving, fixed-factor,
and infinite symbol is $+1$, and the exponent lemma freezes the resulting
progression.  The factor-by-factor identity
\[
 \prod_{p\mid A}(x,d)_p\,\leg{A}{q_1}
   =\leg{-2\varepsilon}{A}\leg{w}{A}
\]
shows that for odd $a$ the break criterion is $\leg{w}{A}=-1$.
Putting $a=7$ gives $A=197$ and the stated fixed-parameter criterion.
\end{proof}
% src: data/l19_tauzero.jsonl (algebraic-criterion and uniform-proof rows)

\begin{proposition}[Finite member closures; PROVED per row]
The non-canonical $\tau=0$ family closes the obstructed cells
$w=131,139,151$ at $a=7$, $q_1=Q=41$, $k=0$.  It also closes
$[179,[-1,1]]$ at $a=7$, $\varepsilon=-1$, $f=179$,
$q_1=Q=251$, $k=0$, with prime rung, tied status exact, and empty
ramified set.
\end{proposition}
% src: data/l18_route131.jsonl; data/l18_horizon_sweep_a.jsonl
% src: data/l19_cell179.jsonl (verified closure row)

These member rows are stronger than the class theorem only at their four
individual cells.  Their fixed factor $f=w$ divides $\num(z)$, so they pass
through the controlled non-coprimality door excluded by the generalized
coprime wall.  The kernel verdict uses the same tied-status, ramification,
alignment, and complete cofactor ladder as the canonical rows; refusals in
the surrounding scans are not evidence.
% src: data/l18_route131.jsonl; data/l18_horizon_sweep_a.jsonl
% src: data/l19_cell179.jsonl

\begin{remark}[Composite $w$]
Composite $w$ is not admissible in the architecture: $w$ indexes a
place, and the residue-field construction has no meaning for composite
modulus --- the break occurs already at the field constructor. The
downstream arithmetic does generalize once $w$ is decomposed into
genuine prime places; as an illustration the pseudo-cell
$[21,[-1,1]]$ closes with the escape route at the honest prime $f = 3$
(member $k=0$, $Q=11$, factorint rung, tied and ramified checks exact).
This is a clarification of scope, not an extension of the horizon.
% src: data/l17_composite_w.jsonl
\end{remark}
